\section{Theory Stack for ELARA-U: Shift-Aware Reliability Routing}
\label{sec:theory}

\subsection{Notation and Setup}

Let $\mathcal{D}=\{1,\ldots,M\}$ be a finite detector zoo. For a task $\tau$, detector $m\in\mathcal{D}$ produces an anomaly score $s_m(x)\in \mathbb{R}$. Let $V_\tau$ denote the validation distribution and $T_\tau$ denote the deployment/test distribution. Let
\[
A_V(m)=\mathrm{AUROC}_{V_\tau}(s_m), \qquad
A_T(m)=\mathrm{AUROC}_{T_\tau}(s_m).
\]
The validation auto-selector is
\[
m_{\mathrm{val}} = \arg\max_{m\in\mathcal{D}} A_V(m),
\]
and the test oracle is
\[
m^\star = \arg\max_{m\in\mathcal{D}} A_T(m).
\]
The test regret of a method $h$ is
\[
\mathrm{Reg}_T(h)=A_T(m^\star)-A_T(h).
\]
When $h$ outputs a fused or stacked score $g_h(x)$ instead of a single detector, define
\[
A_T(h)=\mathrm{AUROC}_{T_\tau}(g_h).
\]
ELARA-U observes validation-only reliability features and unlabeled deployment-shift features, including score dispersion, entropy, inter-detector disagreement, calibration error, detector degeneracy flags, and train-validation/test-score drift statistics. It never observes test labels during routing.

\subsection{Main Evidentiary Theorems}

\begin{theorem}[No Universal Dominance Without Assumptions]
\label{thm:universal_dominance}
For any deterministic detector, fusion rule, or meta-router $h$ that selects or combines scores from a finite detector zoo $\mathcal{D}$, there exists a pair of anomaly-detection tasks $(\tau_1,\tau_2)$ such that $h$ cannot be strictly optimal on both tasks. Therefore, no universal anomaly detector or router can dominate all alternatives across all possible domains without assumptions on task structure, reliability, or shift.
\end{theorem}
\begin{proof}
Fix any method $h$. Since $h$ is deterministic, for a given task representation it outputs either one detector or a fixed scoring function $g_h$ built from detector scores.

Construct a task $\tau_1$ where $g_h$ ranks all anomalous examples above all normal examples. Then
\[
A_{\tau_1}(h)=1.
\]
Now construct a second task $\tau_2$ on the same score space by reversing the anomaly labels: all examples ranked high by $g_h$ are normal and all examples ranked low by $g_h$ are anomalous. Under this label reversal,
\[
A_{\tau_2}(h)=0.
\]
However, there exists another detector $m'$ that ranks the $\tau_2$ anomalies above normals, achieving
\[
A_{\tau_2}(m')=1.
\]
\end{proof}
Thus $h$ is optimal on $\tau_1$ but maximally suboptimal on $\tau_2$. Since this construction can be made for any fixed $h$, no detector, fusion rule, or router can be universally dominant without additional assumptions.

\medskip
\noindent\textbf{Interpretation.}
This theorem justifies the ELARA-U claim structure: the paper does not claim ``best on every dataset.'' The defensible claim is reduction of regret and negative transfer under specified reliability and shift conditions.

\begin{theorem}[Validation Auto-Selection Is Near-Oracle Under Stable Transfer]
\label{thm:stable_transfer}
Assume validation and deployment detector performances are uniformly close:
\[
\sup_{m\in\mathcal{D}} |A_T(m)-A_V(m)| \leq \varepsilon.
\]
Then validation-AUROC auto-selection has deployment regret bounded by
\[
\mathrm{Reg}_T(m_{\mathrm{val}}) \leq 2\varepsilon.
\]
\end{theorem}
\begin{proof}
Let $m^\star=\arg\max_{m}A_T(m)$ be the test oracle. By definition of $m_{\mathrm{val}}$,
\[
A_V(m_{\mathrm{val}})\geq A_V(m^\star).
\]
Using the uniform stability assumption,
\[
A_T(m_{\mathrm{val}})\geq A_V(m_{\mathrm{val}})-\varepsilon,
\]
and
\[
A_V(m^\star)\geq A_T(m^\star)-\varepsilon.
\]
Combining these inequalities gives
\[
A_T(m_{\mathrm{val}})
\geq A_V(m_{\mathrm{val}})-\varepsilon
\geq A_V(m^\star)-\varepsilon
\geq A_T(m^\star)-2\varepsilon.
\]
Therefore,
\[
\mathrm{Reg}_T(m_{\mathrm{val}}) = A_T(m^\star)-A_T(m_{\mathrm{val}}) \leq 2\varepsilon.
\]
\end{proof}

\medskip
\noindent\textbf{Interpretation.}
This theorem explains why validation-AUROC auto-selection is hard to beat on clean i.i.d. validation-to-test transfer. If validation and deployment rankings match, auto-select is already near-oracle.

\begin{theorem}[Stale Validation Causes Auto-Selection Failure Under Shift]
\label{thm:stale_validation}
There exist shifted deployment settings where validation auto-selection incurs arbitrarily large regret up to the AUROC range. In particular, for any $\gamma>0$ and any $\Delta>0$ such that the AUROC values remain in $[0,1]$, there exists a two-detector task with
\[
A_V(1)=a+\gamma,\qquad A_V(2)=a,
\]
but
\[
A_T(1)=a-\Delta,\qquad A_T(2)=a.
\]
Then validation auto-selection chooses detector 1, while the deployment oracle chooses detector 2, and the regret is
\[
\mathrm{Reg}_T(m_{\mathrm{val}})=\Delta.
\]
\end{theorem}
\begin{proof}
Because $A_V(1)=a+\gamma>A_V(2)=a$, validation auto-selection chooses $m_{\mathrm{val}}=1$. At deployment time, however, $A_T(2)=a>A_T(1)=a-\Delta$, so the test oracle is $m^\star=2$. Therefore,
\[
\mathrm{Reg}_T(m_{\mathrm{val}}) = A_T(m^\star)-A_T(m_{\mathrm{val}}) = A_T(2)-A_T(1) = a-(a-\Delta) = \Delta.
\]
Since $\Delta$ can be made large subject only to the AUROC range $[0,1]$, stale validation can cause large auto-selection regret.
\end{proof}

\medskip
\noindent\textbf{Interpretation.}
This theorem formalizes the main ELARA-U motivation: validation auto-selection can be near-oracle when validation and test match, but it can fail badly when deployment shift changes the detector ranking.

\begin{theorem}[Reliability-Separation Improvement Over Stale Auto-Selection]
\label{thm:reliability_separation}
Suppose ELARA-U constructs a reliability predictor $\widehat{A}_T(m)$ for deployment performance using validation-only reliability features and unlabeled deployment-shift features. Assume the predictor is uniformly accurate:
\[
\sup_{m\in\mathcal{D}} |\widehat{A}_T(m)-A_T(m)| \leq \eta.
\]
Let $m_R=\arg\max_{m\in\mathcal{D}}\widehat{A}_T(m)$ be the reliability-routed choice. Then
\[
\mathrm{Reg}_T(m_R)\leq 2\eta.
\]
Moreover, if stale validation auto-selection has regret
\[
\mathrm{Reg}_T(m_{\mathrm{val}}) > 2\eta,
\]
then the reliability router strictly improves over validation auto-selection by at least
\[
\mathrm{Reg}_T(m_{\mathrm{val}})-2\eta.
\]
\end{theorem}
\begin{proof}
Let $m^\star=\arg\max_m A_T(m)$ be the test oracle. Since $m_R$ maximizes $\widehat{A}_T(m)$,
\[
\widehat{A}_T(m_R)\geq \widehat{A}_T(m^\star).
\]
Using the uniform prediction-error assumption,
\[
A_T(m_R)\geq \widehat{A}_T(m_R)-\eta,
\]
and
\[
\widehat{A}_T(m^\star)\geq A_T(m^\star)-\eta.
\]
Combining,
\[
A_T(m_R)
\geq \widehat{A}_T(m_R)-\eta
\geq \widehat{A}_T(m^\star)-\eta
\geq A_T(m^\star)-2\eta.
\]
Therefore,
\[
\mathrm{Reg}_T(m_R) = A_T(m^\star)-A_T(m_R) \leq 2\eta.
\]
If $\mathrm{Reg}_T(m_{\mathrm{val}})>2\eta$, then
\[
\mathrm{Reg}_T(m_{\mathrm{val}})-\mathrm{Reg}_T(m_R) \geq \mathrm{Reg}_T(m_{\mathrm{val}})-2\eta > 0.
\]
Thus the reliability router strictly improves over stale validation auto-selection.
\end{proof}

\medskip
\noindent\textbf{Interpretation.}
This is the main positive theorem. ELARA-U can beat auto-select only when its reliability and drift features predict shifted deployment performance better than stale validation AUROC does.

\begin{theorem}[Safe Switching and Fallback Bound]
\label{thm:safe_switching}
Let $\Delta_{\mathrm{fired}} = A_T(h_{\mathrm{route}})-A_T(h_{\mathrm{default}})$ be the performance gain of routing over a default method on the subset of tasks or inputs where the router fires. Suppose ELARA-U switches away from the default only when a valid one-sided $(1-\alpha)$-confidence lower bound satisfies
\[
\mathrm{LCB}_{1-\alpha}(\Delta_{\mathrm{fired}})>0.
\]
Then, with probability at least $1-\alpha$, every admitted switch has non-negative true improvement:
\[
\Delta_{\mathrm{fired}}>0.
\]
Equivalently, the probability that ELARA-U admits a harmful switch is at most $\alpha$.
\end{theorem}
\begin{proof}
By validity of the one-sided confidence bound,
\[
\Pr\left(\mathrm{LCB}_{1-\alpha}(\Delta_{\mathrm{fired}})\leq \Delta_{\mathrm{fired}}\right)\geq 1-\alpha.
\]
ELARA-U admits a switch only if $\mathrm{LCB}_{1-\alpha}(\Delta_{\mathrm{fired}})>0$. On the event that the confidence bound covers the true value,
\[
\Delta_{\mathrm{fired}} \geq \mathrm{LCB}_{1-\alpha}(\Delta_{\mathrm{fired}}) > 0.
\]
Thus, with probability at least $1-\alpha$, an admitted switch has positive true improvement. Therefore, the probability of admitting a harmful switch is at most $\alpha$.
\end{proof}

\medskip
\noindent\textbf{Interpretation.}
This theorem supports fallback/abstention. ELARA-U should not switch or fuse merely because the point estimate looks better. It should fire only when the benefit is statistically certified.

\begin{corollary}[Operating Boundary of ELARA-U]
\label{cor:boundary}
ELARA-U has no guaranteed advantage over validation-AUROC auto-selection under stable validation-to-test transfer. Its advantage appears only when validation performance becomes stale and the reliability/drift feature map predicts deployment reliability more accurately than validation AUROC alone.
\end{corollary}
\begin{proof}
By Theorem 2, if validation and deployment detector performances are uniformly close, validation auto-selection has regret at most $2\varepsilon$, so little improvement is available.

By Theorem 3, deployment shift can make validation auto-selection suffer regret $\Delta$. By Theorem 4, if ELARA-U predicts shifted deployment reliability with error $\eta$, then its regret is at most $2\eta$. Therefore, ELARA-U improves over auto-selection whenever
\[
\Delta > 2\eta.
\]
Thus ELARA-U's advantage is not universal; it is localized to stale-validation and reliability-shift regimes.
\end{proof}
